\documentclass[10pt,a4paper]{amsart}
\usepackage[margin=19mm]{geometry}
\usepackage{amsmath,amssymb,mathtools}
\usepackage[scr=rsfs,cal=euler]{mathalfa}
\usepackage{xfrac}
\usepackage[unicode,hidelinks,bookmarks=false]{hyperref}
\newcommand{\ie}{i.e.\ }
\newcommand{\cf}{cf.\ }
\newcommand{\resp}{respectively\ }
\newcommand{\Subsection}{\subsection}
\providecommand{\nobreakdash}{\nobreak\hbox{-}}
\setlength{\emergencystretch}{2em}
\multlinegap0pt
\usepackage{mathtools}
\usepackage{amssymb}
\usepackage{xcolor}
\usepackage{tikz}
\usepackage{float}
\newtheorem{lemma}{Lemma}
\title{\Large{\bf A pre-$(n+2)$-angulated category which is not\\[2mm] $(n+2)$-angulated}}
\author{Jian He, Jing He, Panyue Zhou}
\date{Source: arXiv:2608.29221v1 (CC BY 4.0)}
\begin{document}
\maketitle

\noindent\emph{Source and licence.} \Large{\bf A pre-$(n+2)$-angulated category which is not\\[2mm] $(n+2)$-angulated}. Version arXiv:2608.29221v1 (CC BY 4.0), distributed by arXiv under the Creative Commons Attribution 4.0 International licence, http://creativecommons.org/licenses/by/4.0/, as recorded in the arXiv licence metadata for this exact version and shown on the abstract page https://arxiv.org/abs/2608.29221v1. Full licence notice: \texttt{licenses/arxiv-2608.29221-CC-BY-4.0.txt}.\par\smallskip

\noindent\textit{Editorial scope.} Counted unit: the article's lemma lem:odd-twist (source0.tex lines 542--564). The article's complete original proof of it is delivered with it (source0.tex lines 566--597, one \texttt{proof} environment of 32 lines). No mathematics is written, repaired or re-derived: the statement and the proof are the article's own lines, in the article's own notation, and no retained line is substituted or re-flowed.  Retained uncounted as the source-local context the proof needs: lines 484--488.  Nothing was omitted from any retained span, so the package carries no omission record.  No retained line contains a citation, so the wrapper carries no bibliography and no bibliography entry.  No retained line contains a figure environment, an image-inclusion command or drawing code, so no image is delivered and the package declares no asset dependency.  Editorial formatting only, in addition: an AMS \texttt{amsart} preamble that restates the article's own declarations and macros used by the retained lines, namely the environments \texttt{lemma} on separate counters, together with the standard packages those lines need.  The retained environments print in this document as Lemma 1 in order of appearance, whereas the article prints the same items with the numbering of its own full document.  The complete native source of the article as distributed in the arXiv e-print for this exact version is bundled unchanged as \texttt{sources/208835/source0.tex}.
\begin{equation}\label{eq:third-power}
(\Theta_\theta^{[3]})_X
=S^6(\theta_X)\circ S^3(\theta_{\Sigma X})\circ\theta_{\Sigma^2X}\colon
\Sigma^3X\xrightarrow{\sim}S^9X.
\end{equation}

\begin{lemma}\label{lem:odd-twist}
Let $\epsilon:\operatorname{Id}\xrightarrow{\sim}\operatorname{Id}$ be a natural automorphism of $\underline{\operatorname{mod}}\mathcal F$ which commutes with $S$ and satisfies $\epsilon^2=1$. Let $\theta_0:\Sigma\xrightarrow{\sim}S^3$ be an isomorphism of triangle functors and set
\[
(\theta_\epsilon)_X=\epsilon_{S^3X}\circ(\theta_0)_X.
\]
Let $\Theta_0^{[3]}$ and $\Theta_\epsilon^{[3]}$ be the natural isomorphisms defined by \eqref{eq:third-power}. Then
\begin{equation}\label{eq:odd-twist}
(\Theta_\epsilon^{[3]})_X
=\epsilon_{S^9X}\circ(\Theta_0^{[3]})_X.
\end{equation}
Equivalently, if
$
(\lambda_*^{[3]})_X
:=\Omega^9\bigl((\Theta_*^{[3]})_X\bigr):
\Omega^9\Sigma^3X\longrightarrow X,
~ *=0,\epsilon,
$
then
\begin{equation}\label{eq:shifted-odd-twist}
(\lambda_\epsilon^{[3]})_X
=\epsilon_X\circ(\lambda_0^{[3]})_X.
\end{equation}
\end{lemma}

\begin{proof}
Since $\epsilon$ commutes with $S$, expanding the definition gives
\[
(\Theta_\epsilon^{[3]})_X
=\epsilon_{S^9X}S^6((\theta_0)_X)\,
\epsilon_{S^6\Sigma X}S^3((\theta_0)_{\Sigma X})\,
\epsilon_{S^3\Sigma^2X}(\theta_0)_{\Sigma^2X}.
\]
For brevity put
\[
f_1=S^6((\theta_0)_X),~
f_2=S^3((\theta_0)_{\Sigma X}),~
f_3=(\theta_0)_{\Sigma^2X}.
\]
Naturality of $\epsilon$ gives
\[
\epsilon_{S^9X}f_1=f_1\epsilon_{S^6\Sigma X},
~
\epsilon_{S^6\Sigma X}f_2=f_2\epsilon_{S^3\Sigma^2X}.
\]
Consequently,
\[
\begin{aligned}
&\epsilon_{S^9X}f_1\epsilon_{S^6\Sigma X}f_2
\epsilon_{S^3\Sigma^2X}f_3\\
&\qquad =f_1f_2\epsilon_{S^3\Sigma^2X}f_3
=f_1\epsilon_{S^6\Sigma X}f_2f_3
=\epsilon_{S^9X}f_1f_2f_3,
\end{aligned}
\]
where the first equality uses $\epsilon^2=1$. This is \eqref{eq:odd-twist}. Applying $\Omega^9$ and using the commutation of $\epsilon$ with $\Omega$ yields \eqref{eq:shifted-odd-twist}.
\end{proof}

\end{document}
